\section{Methodology}
\label{sec:method}

We propose a statistical framework that reformulates knowledge graph completion evaluation as a large-scale multiple hypothesis testing problem. Our framework is model-agnostic: it operates on the output scores of any KG embedding model and provides false discovery rate (FDR) control for link prediction. We first formalize the problem~(\ref{subsec:formulation}), then develop empirical $p$-value construction from embedding scores~(\ref{subsec:pvalue}), introduce FDR control procedures~(\ref{subsec:bh}--\ref{subsec:locfdr}), and finally define the FDR-based evaluation protocol~(\ref{subsec:protocol}).

\subsection{Problem Formulation: KG Completion as Multiple Hypothesis Testing}
\label{subsec:formulation}

Let $\mathcal{G} = (\mathcal{E}, \mathcal{R}, \mathcal{T})$ denote a knowledge graph, where $\mathcal{E}$ is the set of entities, $\mathcal{R}$ the set of relation types, and $\mathcal{T} \subseteq \mathcal{E} \times \mathcal{R} \times \mathcal{E}$ the set of observed triples.
A KG embedding model learns a scoring function $\hat{g}: \mathcal{E} \times \mathcal{R} \times \mathcal{E} \to \mathbb{R}$ that assigns a real-valued plausibility score to each candidate triple.
For a query $(h, r, ?)$, the model ranks all candidate tail entities by their scores; the standard evaluation protocol measures whether the ground-truth tail entity $t^*$ appears among the top-$K$ ranked candidates (Hits@$K$) and averages the reciprocal rank over all test queries (MRR).

We propose to view this problem through the lens of \emph{multiple hypothesis testing}.
Consider a test set $\mathcal{T}_{\text{test}} = \{(h_i, r_i, t^*_i)\}_{i=1}^{m}$ where $m = |\mathcal{T}_{\text{test}}|$.
For each test triple $i$, we formulate a null hypothesis:
\begin{equation}
    H_0^{(i)}: \text{``The triple } (h_i, r_i, t^*_i) \text{ does not exist in the true KG.''}
    \label{eq:null}
\end{equation}
Rejecting $H_0^{(i)}$ corresponds to the statistical decision that the model has ``discovered'' triple $i$---i.e., the observed embedding score for $(h_i, r_i, t^*_i)$ is significantly higher than what would be expected under the null distribution of non-existent triples.

This formulation yields $m$ simultaneous hypothesis tests. Let $m_0$ be the number of true null hypotheses (triples that genuinely do not exist or are unpredictable), and let $\pi_0 = m_0 / m$ denote the proportion of null cases.
Following the standard multiple testing framework~\cite{benjamini1995controlling, storey2002direct}, we define:
\begin{itemize}
    \item \textbf{False Discovery Proportion (FDP):} the fraction of rejected null hypotheses that are actually true.
    \item \textbf{False Discovery Rate (FDR):} the expected value of FDP: $\text{FDR} = \mathbb{E}[\text{FDP}]$.
\end{itemize}

Our goal is to reject as many $H_0^{(i)}$ as possible (i.e., maximize statistical power) while controlling the FDR at a pre-specified level $\alpha$ (typically $\alpha = 0.05$).
This transforms KG completion evaluation from a purely ranking-based paradigm into a statistically principled framework: instead of only asking ``how high does the true tail rank?'' (MRR), we also ask ``how many test triples can we reliably claim to have discovered?''

\subsection{Empirical $p$-value Construction from Embedding Scores}
\label{subsec:pvalue}

The first technical challenge is that KG embedding scores are not probabilities---they are uncalibrated real numbers whose scale and distribution vary across models.
We must convert each raw score into a valid $p$-value before applying any FDR control procedure.

For each test triple $(h_i, r_i, t^*_i)$, we construct an empirical null distribution by sampling $K$ negative tail entities.
Specifically, we:
\begin{enumerate}
    \item Sample $K$ candidate tail entities $\{t_1^{(i)}, \dots, t_K^{(i)}\}$ uniformly from $\mathcal{E}$, excluding the ground-truth tail $t^*_i$ and all tails known to form valid triples with $(h_i, r_i)$ in the training set (filtered setting).
    \item Compute the model score for each negative triple: $s_{i,k} = \hat{g}(h_i, r_i, t^{(i)}_k)$, for $k = 1, \dots, K$.
    \item Let $s_i^* = \hat{g}(h_i, r_i, t^*_i)$ denote the score of the true triple.
    \item Define the empirical $p$-value as:
    \begin{equation}
        \hat{p}_i = \frac{|\{k : s_{i,k} \geq s_i^*\}| + 1}{K + 1}.
        \label{eq:pvalue}
    \end{equation}
\end{enumerate}

The $+1$ correction in both numerator and denominator ensures that (i) $\hat{p}_i > 0$ for all $i$, avoiding degenerate cases in subsequent FDR procedures, and (ii) the estimator is conservative, i.e., $\mathbb{E}[\hat{p}_i] \geq \Pr_{H_0}(S \geq s_i^*)$, providing valid (though slightly conservative) inference under the null~\cite{davison1997bootstrap}.

The choice $K = 100$ balances precision against computational cost. With $K = 100$, the smallest attainable $p$-value is $1/(K+1) \approx 0.0099$, which provides sufficient granularity for FDR control at conventional $\alpha$ levels ($0.01$--$0.20$).
On an Apple M1 CPU (16\,GB RAM), the full FDR pipeline completes in approximately 3--5 minutes per model on WN18RR ($m = 2{,}924$) and 25--35 minutes on FB15k-237 ($m = 20{,}438$).
For comparison, model training requires 15--25 minutes (20 epochs), so the FDR pipeline adds roughly 15--20\% overhead relative to training---modest given that it requires no model modification or retraining.
This represents roughly 15--20\% overhead relative to training, which is modest given that the FDR framework requires no model modification or retraining.

\subsection{False Discovery Rate Control}
\label{subsec:bh}

\subsubsection{Benjamini-Hochberg Procedure}
\label{subsec:bh-procedure}

The Benjamini-Hochberg (BH) procedure~\cite{benjamini1995controlling} provides finite-sample FDR control under mild dependence assumptions.
Given $m$ $p$-values $\{\hat{p}_1, \dots, \hat{p}_m\}$ and a target FDR level $\alpha$, the BH procedure operates as follows:
\begin{enumerate}
    \item Sort the $p$-values in ascending order: $\hat{p}_{(1)} \leq \hat{p}_{(2)} \leq \cdots \leq \hat{p}_{(m)}$.
    \item Find the largest index $k$ such that $\hat{p}_{(k)} \leq \frac{k}{m} \cdot \alpha$.
    \item Reject all $H_0^{(i)}$ corresponding to $\hat{p}_{(1)}, \dots, \hat{p}_{(k)}$.
\end{enumerate}

The BH theorem guarantees that $\text{FDR} \leq \frac{m_0}{m} \cdot \alpha \leq \alpha$ when the $p$-values are independent or exhibit positive regression dependency on the subset of true null hypotheses (PRDS)~\cite{benjamini2001control}.
Although the negative tail entities are sampled independently for each test triple, the resulting $p$-values share dependence through the trained entity embeddings: the embedding vector of an entity affects scores for all test triples in which it appears.
If two test triples share entities, changes to the shared embedding affect both scores in the same direction, creating positive dependence consistent with the PRDS condition.
A formal verification for general KG embedding architectures is deferred to future theoretical work.

The number of rejected hypotheses, $R(\alpha) = |\{i: \hat{p}_i \leq \hat{p}_{(k)}\}|$, serves as a measure of the model's \emph{statistical power}---the number of test triples that can be reliably distinguished from noise at FDR level $\alpha$.

\subsubsection{Storey's $q$-value}
\label{subsec:qvalue}

While the BH procedure provides a binary reject/accept decision at a fixed $\alpha$, Storey's $q$-value~\cite{storey2002direct, storey2003statistical} offers a more fine-grained measure.
The $q$-value of a hypothesis $H_0^{(i)}$ is defined as the minimum FDR at which $H_0^{(i)}$ would be rejected:
\begin{equation}
    q(\hat{p}_i) = \min_{\alpha \geq \hat{p}_i} \text{FDR}(\alpha).
\end{equation}

A key component of Storey's method is the estimation of $\pi_0 = m_0 / m$, the proportion of true null hypotheses.
For a tuning parameter $\lambda \in (0, 1)$, an unbiased estimator is:
\begin{equation}
    \hat{\pi}_0(\lambda) = \frac{|\{i : \hat{p}_i > \lambda\}|}{m \cdot (1 - \lambda)}.
    \label{eq:pi0}
\end{equation}
The final $\hat{\pi}_0$ is obtained by smoothing over a grid of $\lambda$ values using a cubic spline and taking the limit as $\lambda \to 1$~\cite{storey2002direct}.
The $q$-value for the $i$-th hypothesis is then computed as:
\begin{equation}
    q_{(m)} = \hat{\pi}_0 \cdot \hat{p}_{(m)}, \qquad
    q_{(i)} = \min\left(\hat{\pi}_0 \cdot \frac{m}{i} \cdot \hat{p}_{(i)}, \; q_{(i+1)}\right), \quad i = m-1, \dots, 1.
\end{equation}

In our KG completion context, $\hat{\pi}_0$ has a direct interpretation: it estimates the proportion of test triples for which the embedding model provides no signal above the background noise level.
A small $\hat{\pi}_0$ indicates that the model captures genuine knowledge for most test triples; a large $\hat{\pi}_0$ reveals that a substantial fraction of the test set is indistinguishable from random noise.

\subsection{Empirical Bayes Local False Discovery Rate}
\label{subsec:locfdr}

While $\hat{\pi}_0$ and $q$-values provide global and per-hypothesis FDR control, Efron's empirical Bayes local FDR~\cite{efron2004large, efron2007size} offers a complementary perspective by estimating the posterior probability that a specific hypothesis is null \emph{given its observed test statistic}.

We first transform each $p$-value to a $z$-score using the inverse normal cumulative distribution function:
\begin{equation}
    z_i = \Phi^{-1}(1 - \hat{p}_i),
    \label{eq:zscore}
\end{equation}
where $\Phi^{-1}$ is the standard normal quantile function. Under the null hypothesis $H_0^{(i)}$, we have $z_i \sim \mathcal{N}(0, 1)$; under the alternative, the $z$-scores are shifted toward positive values.

The mixture density of $z$-scores can be expressed as:
\begin{equation}
    f(z) = \pi_0 \cdot f_0(z) + (1 - \pi_0) \cdot f_1(z),
    \label{eq:mixture}
\end{equation}
where $f_0(z) = \phi(z)$ is the standard normal density (null component) and $f_1(z)$ is the alternative density.
The local FDR is then defined as the posterior probability:
\begin{equation}
    \text{fdr}(z) = \frac{\pi_0 \cdot f_0(z)}{f(z)}.
    \label{eq:locfdr}
\end{equation}

In practice, we estimate $f(z)$ via kernel density estimation on the observed $\{z_i\}_{i=1}^m$, compute $\hat{\pi}_0$ using Storey's method, and evaluate $\text{fdr}(z_i)$ at each observed $z$-score.
A triple with $\text{fdr}(z_i) \leq 0.05$ has at most a 5\% posterior probability of being a false discovery---a calibration guarantee that point estimates like MRR cannot provide.

The three FDR measures are complementary:
\begin{itemize}
    \item \textbf{BH rejection count} $R(\alpha)$: how many discoveries at a guaranteed FDR level.
    \item \textbf{Storey's $\hat{\pi}_0$}: what fraction of test triples are noise---a model-level diagnostic.
    \item \textbf{Local FDR}: per-triple reliability---which specific triples are confidently predicted.
\end{itemize}

\subsection{FDR-Based KG Completion Evaluation Protocol}
\label{subsec:protocol}

Building on the above statistical machinery, we define a three-level evaluation protocol that complements (not replaces) the standard MRR/Hits@$K$ metrics.

\subsubsection{Level 1: Global Diagnostic}
For a given KG embedding model on a test set, compute:
\begin{itemize}
    \item $\hat{\pi}_0$ (Eq.~\ref{eq:pi0}): the estimated proportion of unpredictable test triples.
    \item $R(\alpha)$ for $\alpha \in \{0.01, 0.05, 0.10, 0.20\}$: the number of statistically significant discoveries at each FDR level.
\end{itemize}

These metrics reveal the overall reliability of the model's predictions.
For illustration, on WN18RR with TransE, we obtain $\hat{\pi}_0 = 0.295$, meaning approximately 29.5\% of the 2,924 test triples are indistinguishable from noise.
At $\alpha = 0.05$, only $R(0.05) = 262$ triples (9.0\%) are rejected despite an MRR of 0.560.
This notable gap---a model that ranks well yet fails to reliably discover most triples---illustrates the value of FDR-based evaluation.

\subsubsection{Level 2: Relation-Stratified Analysis}
KG relations vary widely in their predictability. We propose to compute $\hat{\pi}_0^{(r)}$ and $R^{(r)}(\alpha)$ for each relation type $r \in \mathcal{R}$ separately.
This stratified analysis reveals which relations the model ``truly understands'' versus those where it performs near random.

Continuing the WN18RR example: for TransE, two relations with similar test-set sizes exhibit dramatically different FDR profiles:
\begin{itemize}
    \item Relation~1 ($n=1074$): $\hat{\pi}_0 = 0.061$, FDR rejection rate 92.7\%---the model reliably discovers nearly all triples.
    \item Relation~3 ($n=1067$): $\hat{\pi}_0 = 0.388$, FDR rejection rate 27.6\%---less than a third of triples are statistically significant, despite comprising 36.5\% of the test set.
\end{itemize}
Traditional MRR would average these two relations into a single number, completely masking this extreme heterogeneity.

\subsubsection{Level 3: Cross-Model FDR Comparison}
For two (or more) KG embedding models, we compare their FDR profiles at each level.
This reveals whether the standard MRR-based model ranking aligns with the FDR-based reliability ranking.
If Model~A has a higher MRR but Model~B rejects more hypotheses at the same FDR level, this signals that MRR alone provides an incomplete picture of model quality.

\subsection{Relationship to Existing Evaluation}
\label{subsec:relationship}

Our FDR-based protocol does not replace MRR or Hits@$K$; rather, it augments them with a complementary dimension of information.
MRR and Hits@$K$ measure \emph{ranking quality}: when the model makes predictions, how often is the correct answer near the top?
Our FDR framework measures \emph{statistical reliability}: for which test triples, and for what proportion of the test set, can we confidently claim the model has made a discovery beyond random chance?

The two perspectives are complementary and together provide a richer characterization:

\begin{center}
\begin{tabular}{p{0.45\textwidth} p{0.45\textwidth}}
\toprule
\textbf{MRR / Hits@$K$ (Traditional)} & \textbf{FDR Framework (This Work)} \\
\midrule
Measures average rank of correct answers & Measures statistical significance of discoveries \\
Single scalar per model & Multi-level: global + per-relation + per-triple \\
Insensitive to per-relation variation & Reveals extreme heterogeneity across relations \\
Ranking-based; no statistical guarantees & Provides FDR control and posterior error probabilities \\
\bottomrule
\end{tabular}
\end{center}

A complete KG completion evaluation should report both---the ranking metrics that have driven progress in the field, and the FDR metrics that reveal \emph{how reliable} those rankings truly are.

\begin{algorithm}[t]
\caption{FDR-Based KG Completion Evaluation Protocol}
\label{alg:fdr-kg}
\begin{algorithmic}[1]
\Require Trained KG embedding model with scoring function $\hat{g}$;
         Test set $\mathcal{T}_{\text{test}} = \{(h_i, r_i, t^*_i)\}_{i=1}^m$;
         Entity set $\mathcal{E}$; number of negative samples $K$; FDR level $\alpha$
\Ensure $p$-values $\{\hat{p}_i\}_{i=1}^m$; BH rejections $R(\alpha)$; Storey $\hat{\pi}_0$, $q$-values; local FDR

\State \textbf{// Phase 1: $p$-value construction}
\For{$i = 1$ \textbf{to} $m$}
    \State Sample $K$ negative tails $\{t_k^{(i)}\}_{k=1}^K \subset \mathcal{E}$ (filtered)
    \State Compute $s_i^* = \hat{g}(h_i, r_i, t^*_i)$ \Comment{True triple score}
    \For{$k = 1$ \textbf{to} $K$}
        \State Compute $s_{i,k} = \hat{g}(h_i, r_i, t^{(i)}_k)$ \Comment{Negative triple score}
    \EndFor
    \State $\hat{p}_i \gets \frac{|\{k: s_{i,k} \geq s_i^*\}| + 1}{K + 1}$ \Comment{Eq.~\ref{eq:pvalue}}
\EndFor

\State \textbf{// Phase 2: Benjamini-Hochberg FDR control}
\State Sort $\{\hat{p}_i\}_{i=1}^m$ ascending: $\hat{p}_{(1)} \leq \cdots \leq \hat{p}_{(m)}$
\State $k^* \gets \max\{k: \hat{p}_{(k)} \leq \frac{k}{m} \cdot \alpha\}$
\State $R(\alpha) \gets k^*$ \Comment{Number of discoveries at FDR level $\alpha$}

\State \textbf{// Phase 3: Storey's $q$-value and $\hat{\pi}_0$ estimation}
\For{$\lambda \in \Lambda$} \Comment{Grid of tuning parameters, e.g., $\Lambda = \{0.1, \dots, 0.9\}$}
    \State $\hat{\pi}_0(\lambda) \gets \frac{|\{i: \hat{p}_i > \lambda\}|}{m \cdot (1-\lambda)}$ \Comment{Eq.~\ref{eq:pi0}}
\EndFor
\State $\hat{\pi}_0 \gets \text{smooth}\{\hat{\pi}_0(\lambda)\}_{\lambda \to 1}$ \Comment{Cubic spline extrapolation}
\State Compute $q$-values from $\{\hat{p}_i\}$ and $\hat{\pi}_0$

\State \textbf{// Phase 4: Efron's local FDR}
\State $z_i \gets \Phi^{-1}(1 - \hat{p}_i)$ for all $i$ \Comment{Eq.~\ref{eq:zscore}}
\State Estimate mixture density $f(z)$ via kernel density estimation
\State $\text{fdr}(z_i) \gets \frac{\hat{\pi}_0 \cdot \phi(z_i)}{f(z_i)}$ for all $i$ \Comment{Eq.~\ref{eq:locfdr}}

\State \textbf{return} $\{\hat{p}_i\}, R(\alpha), \hat{\pi}_0, \{q_i\}, \{\text{fdr}(z_i)\}$
\end{algorithmic}
\end{algorithm}
